\documentclass[leqno]{article}
% Shared mark scheme preamble for resources.danieljsmith.org: the pack
% preamble, unchanged, then the scheme's own layer. A scheme's .tex is
% rendered from its YAML (scripts/README.md, Mark schemes), never edited.
\input{\CurrentFilePath/pack_preamble.tex}

% The page: equal side margins, so the table runs the full width of the
% header rule. The footers' right edge follows the text's.
\geometry{left=0.8in, right=0.8in}
\makeatletter
\renewcommand{\djs@footR}[1]{%
  \put(\dimexpr1in+\hoffset+\oddsidemargin+\textwidth\relax,
    -\dimexpr\paperheight-0.32in\relax){\makebox[0pt][r]{#1}}}
\makeatother

\usepackage[export]{adjustbox}
\usepackage{tabularray}
\definecolor{msGrid}{RGB}{176, 172, 164}
\definecolor{msTint}{RGB}{245, 243, 239}

% The header on every page, the front page included, and the booklet label
% under the contents-page title. \djsQuestionPageStyle stays empty.
\newcommand{\djsSchemeHeader}[1]{%
  \djsPackHeader[markscheme]{\textit{Mark scheme} - #1}%
  \renewcommand{\djsPackTitle}{#1}%
  \renewcommand{\djsBookletLabel}{Mark scheme}%
}

% The front page: the contents heading, spaced as tocloft spaces the packs'
% "Questions", then the contents table. A question's name and page number
% both link to its first page.
\newcommand{\msContentsHeading}{%
  \vskip\cftbeforetoctitleskip
  {\noindent\cfttoctitlefont\contentsname\par}%
  \vskip\cftaftertoctitleskip}
\NewTblrEnviron{mscontents}
\SetTblrInner[mscontents]{
  colspec = {Q[wd=26mm, l] X[l] Q[wd=16mm, c] Q[wd=14mm, r]},
  width = \linewidth,
  rowsep = 4pt, colsep = 6pt,
  hlines = {0.4pt, msGrid}, vlines = {0.4pt, msGrid},
  hline{2} = {2.2pt, scarlet},
  row{1} = {font=\small\bfseries},
  row{Z} = {font=\bfseries},
  column{2} = {font=\small},
}
\newcommand{\msContentsQuestion}[1]{\hyperref[q:#1]{Question #1}}
\newcommand{\msContentsPage}[1]{\hyperref[q:#1]{\pageref*{q:#1}}}
\newcommand{\msSep}{\hspace{0.9em}}

% The conventions pages' layout: headings in the documents' own bold, tables
% ruled and padded as the scheme's, roomy lists.
\newcommand{\msHead}[1]{%
  \par\vspace{18pt}%
  {\noindent\bfseries #1\par}%
  \vspace{6pt}}
\NewTblrEnviron{mstblr}
\SetTblrInner[mstblr]{rowsep = 6pt, colsep = 6pt,
  hlines = {0.4pt, msGrid}, vlines = {0.4pt, msGrid}}
\newenvironment{msList}
  {\begin{itemize}[leftmargin=1.2em, itemsep=5pt, topsep=3pt, parsep=0pt]}
  {\end{itemize}}

% A question's first page: the link target and the PDF bookmark, as
% \questionitem gives the packs.
\newcommand{\msQnum}{}
\newcommand{\msQuestion}[1]{%
  \clearpage
  \renewcommand{\msQnum}{#1}%
  \phantomsection\label{q:#1}%
  \addcontentsline{toc}{subsection}{\protect\djsTocEntry{#1}}}

% The scheme's table: Question 16mm, Scheme the rest, Marks 16mm, Notes
% 44mm; every cell ruled in msGrid; the scarlet rule under the head row. A
% long table has no caption or "continued" foot; a continuation page opens
% with a muted line above the repeated head row. The rules across a part's
% rows, the alternative's tint and the total row are the rendered table's.
\NewTblrTheme{ms}{%
  \DefTblrTemplate{firsthead}{default}{}%
  \DefTblrTemplate{middlehead,lasthead}{default}{%
    {\small\itshape\color{djsMuted}Question \msQnum\ continued\par}%
    \vspace{6pt}}%
  \DefTblrTemplate{firstfoot,middlefoot,lastfoot}{default}{}%
}
\NewTblrEnviron{msscheme}
\SetTblrOuter[msscheme]{long, theme=ms, entry=none, label=none}
\SetTblrInner[msscheme]{
  colspec = {Q[wd=16mm, l] X[l] Q[wd=16mm, c] Q[wd=44mm, l]},
  width = \linewidth,
  rowsep = 6pt, colsep = 6pt,
  rowhead = 1,
  vlines = {0.4pt, msGrid},
  hline{1} = {0.4pt, msGrid}, hline{2} = {2.2pt, scarlet},
  hline{Z} = {0.4pt, msGrid},
  column{1} = {font=\bfseries}, column{3} = {font=\bfseries},
  column{4} = {font=\small},
  row{1} = {rowsep=4pt, font=\small\bfseries},
}

% Two marks on a row: each note starts with its code.
\newcommand{\msNote}[2]{\textbf{#1}\enspace #2}
% A full-row figure: the worked solution's own picture, scaled uniformly and
% only down, to the table less the Question column. The picture is drawn
% with the line width it has in the solutions document, at its list depth
% there (the optional argument, 1 by default: a part's working), so one
% sized by \linewidth keeps its size; tabularray sets the cell's \linewidth
% to the cell's measured width, which would shrink it at every pass. The
% widths are measured once, in a solution box at the packs' text width
% (pack_preamble.tex's margins, 0.8in and 1in).
\newlength{\msSpanW}
\makeatletter
\newcommand{\ms@keepLW}[1]{%
  \expandafter\xdef\csname ms@solnLW@#1\endcsname{\the\linewidth}}
\AtBeginDocument{%
  \setlength{\msSpanW}{\dimexpr\textwidth-16mm-28pt\relax}%
  \setbox\z@\vbox{%
    \begin{minipage}{\dimexpr\paperwidth-1.8in\relax}
      \begin{djsSolution}
        \ms@keepLW{0}%
        \begin{enumerate}\item \ms@keepLW{1}%
        \begin{enumerate}\item \ms@keepLW{2}%
        \begin{enumerate}\item \ms@keepLW{3}%
        \end{enumerate}\end{enumerate}\end{enumerate}
      \end{djsSolution}
    \end{minipage}}}
\newcommand{\msFigure}[2][1]{%
  \adjustbox{max width=\msSpanW, max totalheight=0.6\textheight, valign=t}{%
    \hsize=\csname ms@solnLW@#1\endcsname
    \linewidth=\hsize
    #2}}
\makeatother
% A figure mark's feature, in the Scheme column.
\newcommand{\msFeature}[1]{\textit{#1}}
% An alternative's title row.
\newcommand{\msAlternative}[1]{\textit{Alternative:} #1}
% Drafts only: a queried mark's code, the query's why under its note, and a
% flag in its part's place.
\newcommand{\msQueried}[2]{#1~\(\to\) #2?}
\newcommand{\msQuery}[1]{{\itshape\color{djsMuted}Query: #1}}
\newcommand{\msFlag}[2]{\textbf{#1:}\enspace #2}


\djsSchemeHeader{Mark Scheme Fixture}
\djsSitePath{Core Pure}{Polar Coordinates}

\djsLastUpdated{18 October 2026}
\begin{document}
\djsSetQuestionCount{4}
\djsAddMarks{1}{7}
\djsAddMarks{2}{6}
\djsAddMarks{3}{8}
\djsAddMarks{4}{18}
\vspace*{6pt}
\djsContentsTitle
\vspace{30pt}
\msContentsHeading
\begin{mscontents}{}
Question & Marks by part & Marks & Page \\
\msContentsQuestion{1} & (a)\,3\msSep (b)\,4 & 7 & \msContentsPage{1} \\
\msContentsQuestion{2} & (a)\,4\msSep (b)\,2 & 6 & \msContentsPage{2} \\
\msContentsQuestion{3} & & 8 & \msContentsPage{3} \\
\msContentsQuestion{4} & (a)\,16\msSep (b)\,2 & 18 & \msContentsPage{4} \\
\SetCell[c=2]{r} Total & & 39 & \\
\end{mscontents}
\newpage
% The conventions pages, shared by every pack's mark scheme: one source of
% truth, \input after the front page, with no strand or pack logic. The keys
% on the first side, the rules on the second. The scheme layer defines
% \msHead, msList and mstblr.
\vspace*{6pt}%
{\noindent\cfttoctitlefont Reading the mark scheme\par}%
\vspace{10pt}%
These rules apply throughout, unless a note says otherwise.

\msHead{Marks}
Every mark is worth one, so \textbf{M1 A1} on a row is two marks. Only marks
coded ft follow through.\par
\vspace{10pt}%
\begin{mstblr}{
  colspec = {Q[wd=15mm, l] X[l]},
  width = \linewidth,
  hline{2} = {2.2pt, scarlet},
  column{1} = {font=\bfseries},
  row{1} = {font=\small\bfseries, rowsep=4pt},
}
Code & Meaning \\
M1 & Method: a correct method, applied to this question. A slip doesn't cost it if the method is clear. A formula must be used, not just quoted. \\
dM1 & Dependent method: as M1, but only if the method mark before it was earned, unless its note names another. \\
A1 & Accuracy: a correct answer or step. It needs the method mark it follows. \\
A1ft & As A1, following through from an earlier answer that the note names. \\
A1~AG & The last mark of a part whose answer is printed in the question: complete working, with no errors. \\
B1 & A correct value, fact or statement. It needs no method mark. \\
B1ft & As B1, following through from an earlier answer that the note names. \\
E1 & A written explanation, interpretation or conclusion. Concluding a proof by induction needs every earlier mark. \\
\end{mstblr}

\msHead{Abbreviations}
ft and AG qualify a code in the Marks column. The others appear in the
notes.\par
\vspace{10pt}%
\begin{mstblr}{
  colspec = {Q[wd=17mm, l] X[l] Q[wd=15mm, l] X[l]},
  width = \linewidth,
  column{1} = {font=\bfseries}, column{3} = {font=\bfseries},
}
ft & follow through & AG & answer given in the question \\
cao & correct answer only, with no follow-through & cso & correct solution only: no errors in the part \\
oe & or an equivalent form & awrt & any answer that rounds to this value \\
soi & seen, or implied by later working & isw & ignore working after a correct answer \\
SC & special case: credit outside the main method & their & the student's own earlier value \\
condone & a slip that doesn't cost the mark & [\;\,] & working the mark doesn't need \\
s.f. & significant figures & d.p. & decimal places \\
\end{mstblr}

\newpage
\msHead{Reading the scheme}
\begin{msList}
\item Each row is a line of working, with the marks it earns beside it.
  Where a row carries two marks, each note starts with its code.
\item \textbf{[\textit{n}]} ends each part with its total.
\item A shaded \textit{Alternative} is a second method for the same total.
  Mark an answer by the method it follows.
\item Any other valid method earns the equivalent marks, unless the question
  names a method or says \textit{hence}.
\item A sketch earns one B1 for each feature named below its figure.
\end{msList}

\msHead{Accuracy and answers}
\begin{msList}
\item Unless a question says otherwise, give answers to 3 significant
  figures, and angles in degrees to 1 decimal place. Angles in radians take
  3 significant figures.
\item A final answer with more figures is accepted if it rounds to the
  scheme's value; one with fewer is not. An answer spoiled by rounding too
  early loses its accuracy mark.
\item Where a question asks for an exact answer, a decimal loses the final
  accuracy mark. Otherwise either is accepted.
\item Wrong or missing units cost nothing unless the question asks for them.
\item A correct answer with no working earns full marks, unless the question
  needs the working: show that, prove, hence, a method it names, or an exact
  answer. A wrong answer with no working earns nothing. Calculator answers
  are treated the same way.
\item A correct answer, once reached, isn't spoiled by later working (isw).
\end{msList}

\msHead{Value of $\boldsymbol{g}$}
\begin{msList}
\item Take $g = 9.8\,\mathrm{m\,s^{-2}}$ unless a question states a value,
  which is then the only one accepted. Answers that use $g = 9.8$ may be
  given to 2 or 3 significant figures.
\item Where a question doesn't state $g$, an answer worked with $g = 10$
  throughout earns the same marks, keeping any value printed in the
  question. The worked solutions list the final answers with $g = 10$.
\end{msList}


\msQuestion{1}
\begin{msscheme}{
  hline{3-6,8-11} = {2-Z}{0.4pt, msGrid},
  hline{7,12} = {0.4pt, msGrid},
  row{12} = {font=\bfseries},
}
Question & Scheme & Marks & Notes \\
1(a) & \SetCell[c=3]{c} \msFigure{\begin{tikzpicture}[scale=1.5]
\def\ang{120}
\coordinate (O) at (0,0);
\coordinate (P) at (\ang:1);
\coordinate (Q) at (4,0);
\draw[->] (-1.2,0) -- (4.7,0) node[right] {$x$};
\draw[->] (0,-2.6) -- (0,2.6) node[above] {$y$};
\draw[thick,domain=0:360,samples=300,smooth,variable=\t]
plot ({2*(1+cos(\t))*cos(\t)},{2*(1+cos(\t))*sin(\t)});
\draw[thick,->] (O) -- ++(\ang:3) node[above left] {$C_2:\ \theta=\frac{2\pi}{3}$};
\node[below right] at (O) {$O$};
\fill (Q) circle (1.2pt);
\node[below right] at (Q) {$(4,0)$};
\fill (P) circle (1.2pt);
\node[below left] at (P) {$\left(1,\frac{2\pi}{3}\right)$};
\node[above] at (2.0,1.65) {$C_1$};
\end{tikzpicture}} & & \\
& \msFeature{Shape} & B1 & {A cardioid, symmetric about the initial line, with its cusp at the pole.} \\
& \msFeature{Labelled point} & B1 & {\((4, 0)\) marked, or \(r = 4\) on the initial line.} \\
& \msFeature{Half-line} & B1 & {From the pole at angle \(\frac{2\pi}{3}\), meeting \(C_1\) once more.} \\
\nopagebreak & & [3] & \\
1(b) & {\(A = \frac12\int_{2\pi/3}^{\pi} 4(1 + \cos\theta)^2\,\mathrm{d}\theta\)} & M1 & {Uses \(\frac12\int r^2\,\mathrm{d}\theta\) with \(r = 2(1 + \cos\theta)\). Limits may come later.} \\
& {\(\phantom{A} = \int_{2\pi/3}^{\pi} (3 + 4\cos\theta + \cos 2\theta)\,\mathrm{d}\theta\)} & M1 & {Uses \(\cos^2\theta = \frac12(1 + \cos 2\theta)\) to reach an integrable form.} \\
& {\(\phantom{A} = \left[3\theta + 4\sin\theta + \frac12\sin 2\theta\right]_{2\pi/3}^{\pi}\)} & A1 & \\
& {\(\phantom{A} = \pi - \frac{7\sqrt3}{4}\) [\(\approx 0.111\)]} & A1 & {cso, from the limits \(\frac{2\pi}{3}\) and \(\pi\).} \\
\nopagebreak & & [4] & \\
\nopagebreak \SetCell[c=2]{r} Total for Question 1 & & 7 & \\
\end{msscheme}

\msQuestion{2}
\begin{msscheme}{
  hline{3-5,7-8} = {2-Z}{0.4pt, msGrid},
  hline{6,9} = {0.4pt, msGrid},
  row{9} = {font=\bfseries},
}
Question & Scheme & Marks & Notes \\
2(a) & {\(\text{EPE}_0 = \frac{2mg(2a)^2}{2(2a)} = 2mga\); the string is slack at \(\frac32 a\) above \(O\)} & \msQueried{B1}{M1} & {The initial elastic energy, and none once \(P\) is \(\frac32 a\) above \(O\).\\ \msQuery{The initial energy needs Hooke's formula applied, a method.}} \\
& {\(4mga + 2mga = \frac12 mv^2 + \frac32 mga\)} & M1 \msQueried{A1}{A1ft} & {\msNote{M1}{Forms an energy equation with kinetic, gravitational and elastic terms, each of the right form.}\\ \msNote{A1}{\msQuery{The equation may carry an earlier wrong elastic energy.}}} \\
& {\(v = 3\sqrt{ag}\)} & A1 & \\
\nopagebreak & & [4] & \\
2(b) & {\(T = \frac{2mg(2a)}{2a} = 2mg\)\\[4pt] \(m\alpha = 2mg + mg\)} & M1 & {Uses Hooke's law for the tension and applies Newton's second law with the tension and the weight.} \\
& {\(\alpha = 3g\), downwards} & A1 & \\
\nopagebreak & & [2] & \\
\nopagebreak \SetCell[c=2]{r} Total for Question 2 & & 6 & \\
\end{msscheme}

\msQuestion{3}
\begin{msscheme}{
  hline{3-15} = {2-Z}{0.4pt, msGrid},
  hline{16} = {0.4pt, msGrid},
  cell{9-15}{2-4} = {bg=msTint},
  row{16} = {font=\bfseries},
}
Question & Scheme & Marks & Notes \\
3 & {\(\overrightarrow{CD} = 2\mathbf i + \mathbf j - \mathbf k\), \(\overrightarrow{CE} = \mathbf i + 3\mathbf j\)} & M1 & {Finds two direction vectors in the plane by subtracting position vectors of the three points.} \\
& {\(\mathbf n = \overrightarrow{CD}\times\overrightarrow{CE} = 3\mathbf i - \mathbf j + 5\mathbf k\)} & M1 A1 & {\msNote{M1}{Finds the vector product of their two direction vectors, or solves for a normal by scalar products.}} \\
& {\(\mathbf r\cdot(3\mathbf i - \mathbf j + 5\mathbf k) = 14\), or \(3x - y + 5z = 14\)} & A1 & \\
& {\(\mathbf p' = \mathbf p - 2\,\dfrac{\mathbf p\cdot\mathbf n - d}{\mathbf n\cdot\mathbf n}\,\mathbf n\)} & M1 & {Uses the reflection formula with their plane, or another full method for the image of \(A\).} \\
& {\(\mathbf p\cdot\mathbf n - d = 35\), \(\mathbf n\cdot\mathbf n = 35\)} & A1 & \\
& {\(\overrightarrow{OB} = (4\mathbf i - 2\mathbf j + 7\mathbf k) - 2(3\mathbf i - \mathbf j + 5\mathbf k)\)} & dM1 & {Substitutes their values into the formula, with the factor of 2 and the sign right.} \\
& {\(B = (-2, 0, -3)\)} & A1 & \\
& \SetCell[c=3]{l} \msAlternative{the line through \(A\) along the normal} & & \\
\nopagebreak & {\(\mathbf n = \overrightarrow{CD}\times\overrightarrow{CE} = 3\mathbf i - \mathbf j + 5\mathbf k\)} & M1 M1 & {\msNote{M1}{Finds two direction vectors in the plane from the three points.}\\ \msNote{M1}{Finds their vector product, or solves for a normal by scalar products.}} \\
& {\(\mathbf n = 3\mathbf i - \mathbf j + 5\mathbf k\); \(3x - y + 5z = 14\)} & A1 A1 & \\
& {\(\mathbf r = (4\mathbf i - 2\mathbf j + 7\mathbf k) + \lambda(3\mathbf i - \mathbf j + 5\mathbf k)\)\\[4pt] \(3(4 + 3\lambda) - (-2 - \lambda) + 5(7 + 5\lambda) = 14\)} & M1 & {Forms the line through \(A\) along their normal and substitutes it into their plane.} \\
& {\(\lambda = -1\): the line meets the plane at \(C(1, -1, 2)\)} & A1 & \\
& {\(\overrightarrow{OB} = \overrightarrow{OA} + 2\overrightarrow{AC}\)} & dM1 & {Finds \(B\) by doubling the step from \(A\) to the plane, or uses \(C\) as the midpoint of \(AB\).} \\
& {\(B = (-2, 0, -3)\)} & A1 & \\
\nopagebreak \SetCell[c=2]{r} Total for Question 3 & & 8 & \\
\end{msscheme}

\msQuestion{4}
\begin{msscheme}{
  hline{3-18,20} = {2-Z}{0.4pt, msGrid},
  hline{19,21} = {0.4pt, msGrid},
  row{21} = {font=\bfseries},
}
Question & Scheme & Marks & Notes \\
4(a) & {Along the line of centres, \(\frac{1}{\sqrt2}(\mathbf i + \mathbf j)\):\\[4pt] \(P\): \(\frac{6 + 2}{\sqrt2}\), \(Q\): \(\frac{-2 + 2}{\sqrt2}\)} & M1 & {Resolves both velocities along the line of centres, using a unit vector parallel to \(\mathbf i + \mathbf j\).} \\
& {\(P\): \(4\sqrt2\), \(Q\): \(0\)} & A1 & {Both parallel components, in \(\mathrm{m\,s^{-1}}\).} \\
& {Perpendicular to it, \(\frac{1}{\sqrt2}(-\mathbf i + \mathbf j)\):\\[4pt] \(P\): \(\frac{-6 + 2}{\sqrt2}\), \(Q\): \(\frac{2 + 2}{\sqrt2}\)} & M1 & {Resolves both velocities perpendicular to the line of centres.} \\
& {\(P\): \(-2\sqrt2\), \(Q\): \(2\sqrt2\)} & A1 & {Both perpendicular components, with signs consistent with their unit vector.} \\
& {The spheres are smooth, so the impulse is along the line of centres: each perpendicular component is unchanged} & E1 & {Explains that smooth spheres give no impulse perpendicular to the line of centres, so those components are unchanged.} \\
& {\(0.4(4\sqrt2) + 0.6(0) = 0.4v_P + 0.6v_Q\)} & M1 & {Uses conservation of momentum along the line of centres, with their parallel components. Condone a sign slip.} \\
& {\(8\sqrt2 = 2v_P + 3v_Q\)} & A1 & \\
& {\(v_Q - v_P = 4\sqrt2\,e\)} & M1 & {Uses Newton's law of restitution along the line of centres, separation speed over approach speed.} \\
& {\(\frac12(0.4)(4\sqrt2)^2 - \frac12(0.4)v_P^2 - \frac12(0.6)v_Q^2 = 2.88\)} & M1 & {Forms an energy equation in the parallel components, or in the full speeds, with the loss of \(2.88\,\mathrm{J}\).} \\
& {\(v_P = \frac{4\sqrt2}{5}(2 - 3e)\), \(v_Q = \frac{8\sqrt2}{5}(1 + e)\)} & dM1 & {Solves the momentum and restitution equations and substitutes into the energy equation, reaching an equation in \(e\) alone.} \\
& {\(2v_P^2 + 3v_Q^2 = 35.2\)\\[4pt] \((2 - 3e)^2 + 6(1 + e)^2 = \frac{55}{4}\)} & A1 & {A correct quadratic in \(e\), in any form.} \\
& {\(e = \frac12\), as \(e \geq 0\)} & A1 & \\
& {\(v_P = \frac{2\sqrt2}{5}\), \(v_Q = \frac{12\sqrt2}{5}\)} & A1ft & {ft their \(e\), in their expressions for \(v_P\) and \(v_Q\), with \(0 \leq e \leq 1\).} \\
& {\(\mathbf v_P = \frac{2\sqrt2}{5}\cdot\frac{1}{\sqrt2}(\mathbf i + \mathbf j) + 2\sqrt2\cdot\frac{1}{\sqrt2}(\mathbf i - \mathbf j)\)} & M1 & {Combines the parallel and perpendicular components of the velocity of \(P\) in terms of \(\mathbf i\) and \(\mathbf j\).} \\
& {\(\mathbf v_P = (2.4\mathbf i - 1.6\mathbf j)\,\mathrm{m\,s^{-1}}\)} & A1 & \\
& {\(\mathbf v_Q = (0.4\mathbf i + 4.4\mathbf j)\,\mathrm{m\,s^{-1}}\)} & A1ft & {ft their \(v_Q\), combined with the unchanged perpendicular component as for \(P\).} \\
\nopagebreak & & [16] & \\
4(b) & \SetCell[c=3]{l} \msFlag{Not drafted}{The drafting run stopped before this part.} & & \\
\nopagebreak & & [2] & \\
\nopagebreak \SetCell[c=2]{r} Total for Question 4 & & 18 & \\
\end{msscheme}

\end{document}
